% ============================================================================
% folha-rosto.tex
% ----------------------------------------------------------------------------
% Gera a FOLHA DE ROSTO: a segunda página do documento, com um breve texto
% de apresentação do relatório e os dados da equipe responsável.
%
% Usa dados definidos em edicao.tex (\nomeProjeto, \nomeCliente,
% \gerenteProjeto, \equipeTecnica) e em config.tex (\entidadeExecutora,
% \dataEmissaoDocumento).
%
% Você normalmente NÃO precisa editar este arquivo, apenas os dados que ele
% utiliza (em config.tex e edicao.tex).
% ============================================================================

\thispagestyle{empty}   % folha de rosto não tem cabeçalho/rodapé

\vspace*{2cm}

% -- Texto de apresentação do relatório ---------------------------------------
\begin{center}
\begin{minipage}{0.85\textwidth}
    \setstretch{1.15}
    \noindent Relatório Técnico Conceitual apresentado como entrega
    formal do projeto \textbf{\nomeProjeto}, desenvolvido para a
    empresa \textbf{\nomeCliente}.
\end{minipage}
\end{center}

\vfill

% -- Dados da equipe responsável -----------------------------------------------
\noindent
\begin{minipage}{\textwidth}
    \textbf{Entidade Executora:} \entidadeExecutora \\[0.3cm]
    \textbf{Gerente de Projeto:} \gerenteProjeto \\[0.3cm]
    \textbf{Equipe Técnica:} \equipeTecnica
\end{minipage}

\vspace{1.5cm}

% -- Data de emissão --------------------------------------------------------
\noindent
\begin{minipage}{\textwidth}
    \textbf{Data de Emissão:} \dataEmissaoDocumento
\end{minipage}

\vspace*{1cm}

\clearpage  % garante que o próximo conteúdo (sumário) comece em página nova